\section{The Autonomous Drift Problem}

\subsection{Definitions and Formal Characterization}

In recursive spawning systems, \textbf{drift} denotes a condition in which an agent
instance loses its causal chain back to a human anchor point. We distinguish two
primitive states:

\begin{itemize}
\item \textbf{Orphaned agent:} an agent $a$ whose parent agent $p$ has either
  terminated normally, crashed, or become unreachable, such that $p$ can no
  longer observe, correct, or terminate $a$.

\item \textbf{Drifted agent:} an agent $a$ for which the transitive closure of
  the parent--child relationship does not reach any human principal $h \in H$,
  where $H$ denotes the set of human users with oversight authority over the
  spawning graph.
\end{itemize}

An agent may be orphaned without being drifted (if a living sibling or
supervisor retains a handle), but any drifted agent is necessarily orphaned at
its immediate parent level or at some higher level in its ancestry.

\begin{definition}[Drift]
Let $G = (V, E)$ be a directed acyclic graph (DAG) where each vertex $v \in V$
represents an agent instance and each directed edge $(p, c) \in E$ represents a
spawning relation from parent $p$ to child $c$. Let $H \subset V$ denote the set
of human-anchored root nodes. For any node $v \in V$, define the ancestry
function $A(v)$ as the set of all nodes reachable by following parent edges
transitively to their origin. Then $v$ is \textbf{drifted} if and only if

\[
A(v) \cap H = \varnothing .
\]

An agent is \textbf{orphaned} if its immediate parent $p$ satisfies
$\text{status}(p) \in \{\text{TERMINATED}, \text{UNREACHABLE}\}$.
\end{definition}

Drift can therefore be detected by traversing $G$ toward its roots; if no human
anchor is encountered before a terminal node, the agent has drifted.

\subsection{Conditions That Cause Drift}

Drift does not require malicious behavior. It emerges from four routine
operational conditions:

\begin{enumerate}
\item \textbf{Parent crash:} The spawning process completes and returns a handle
  to the child, but the parent process terminates before the child can be
  registered with a supervisor or logged to a stable registry. The child
  survives and continues executing, but no living entity holds its reference.

\item \textbf{Asynchronous spawn with fire-and-forget semantics:} The parent
  spawns a child and receives a future or promise, but does not block on the
  child's registration. If the parent is deallocated before registration
  completes, the child's reference is lost even though the child remains active.

\item \textbf{Network partition in distributed spawning:} A supervisor spawns
  agents across multiple nodes. A child on node $B$ reports to a monitor on node
  $C$, but the monitor's connection to the supervisor on node $A$ is severed. The
  child continues running; the supervisor believes it is terminated; no human
  sees the child.

\item \textbf{Supervisor eviction:} A supervisor agent is itself spawned by a
  parent and operates for a bounded window (e.g., a workflow-scoped helper). When
  the window expires, the supervisor is terminated. All children it has not
  independently registered with a higher authority become orphaned.
\end{enumerate}

These conditions are not pathological; they represent standard failure modes in
any production-grade agent runtime. Drift is a structural consequence of these
conditions in any system where child agents can outlive their parents.

\subsection{Failure Modes}

Once an agent is drifted, the failure modes compound:

\begin{description}
\item[Resource leak:] A drifted agent holds memory, compute, and potentially
  external credentials (API keys, file handles) with no entity responsible for
  cleanup. Over time, accumulated drifted agents exhaust system resources.

\item[Zombie actions:] A drifted agent continues executing actions—sending
  emails, posting to APIs, modifying state—using credentials inherited from its
  parent chain. Because no human oversees it, these actions occur without
  authorization, potentially causing legal, financial, or reputational harm.

\item[Orphan chain propagation:] A drifted agent may itself spawn children. These
  grandchildren have no human anchor either, propagating drift downward. A single
  orphaned parent can generate an exponentially growing tree of drifted agents.

\item[Observation gap:] Monitoring systems that track agent health by querying
  parent processes find no parent and may incorrectly conclude the child is also
  terminated. Drift is therefore invisible to standard liveness checks.
\end{description}

The zombie action failure mode is particularly consequential for commercial
deployments. An agent authorized to send notifications on behalf of a user
remains authorized after its parent dies. Nothing signals to the credential
store that the agent should be revoked. The agent continues to act.

\subsection{Why Depth Limits Alone Do Not Solve Drift}

A common architectural response to unbounded spawning is to impose a maximum
spawn depth. If a parent may not spawn beyond depth $D$, the reasoning goes, the
longest possible drift chain has length $D$, and a human can simply audit depth
$D$ agents to bound the blast radius.

This reasoning is incorrect for three reasons.

First, \textbf{depth limits constrain the tree's height, not its width}. A single
drifted agent at depth $D$ that spawns breadth-first can create $N$ children at
depth $D+1$, all of them drifted. Depth limits place no bound on the number of
agents that can exist simultaneously at any given depth. A $D=3$ limit with a
drifted root agent spawning 1,000 children at depth 3 produces 1,000 drifted
agents—well within the limit and invisible to depth-based audits.

Second, \textbf{depth limits do not address temporal drift}. An agent at depth
$1$ that runs for 72 hours and performs a critical action (send \$50,000 wire
transfer) at hour 71 is drifting not in the structural sense but in the
temporal sense: no human has observed its behavior for an extended period. Depth
limits place no bound on the time an agent operates between human observations.

Third, \textbf{depth limits are defined relative to a spawning event, not a
human event}. If a drifted agent at depth $D$ spawns, the new agent is at depth
$D+1$. The human at the root may have a policy that depth $>D$ agents must be
approved. The drifted agent does not enforce this policy. The system has no
mechanism to enforce depth limits on descendants of drifted agents because the
drift itself breaks the enforcement path.

\begin{center}
\begin{tabular}{lcc}
\toprule
Property & Depth Limit Satisfied? & Drift Prevented? \\ \midrule
Maximum ancestry length & Yes (bounded by $D$) & No \\
Number of drifted agents & No (unbounded breadth) & No \\
Temporal observation gap & No (no time bound) & No \\
Enforcement path integrity & No (broken by drift) & No \\
\bottomrule
\end{tabular}
\end{center}

Depth limits reduce the structural \emph{distance} of drift but do not reduce
the \emph{probability}, \emph{extent}, or \emph{consequences} of drift once it
occurs.

\subsection{A Simple Formal Model}

We formalize drift as a property of the spawning graph under a failure model.

Let $G_t = (V_t, E_t)$ denote the spawning graph at time $t$. Each vertex $v$
carries a state $s(v) \in \{\text{ACTIVE}, \text{TERMINATED}, \text{UNREACHABLE}\}$.
Let $P(v) = \{u \in V_t \mid (u, v) \in E_t\}$ denote the set of parents of $v$,
and let $S(v) = \{w \in V_t \mid (v, w) \in E_t\}$ denote the set of direct
children.

\begin{axiom}[Parent Failure]
A parent $p \in P(v)$ becomes UNREACHABLE with probability $p_f$ per time unit,
independent of the child's state.
\end{axiom}

\begin{axiom}[Child Autonomy]
A child $v$ continues executing and may spawn descendants regardless of the
state of any $p \in P(v)$.
\end{axiom}

\begin{axiom}[No Automatic Registration]
There is no guarantee that a child $v$ registers with any entity outside its own
ancestry chain upon parent failure.
\end{axiom}

\begin{definition}[Drift Event]
A drift event occurs at time $t$ for agent $v$ if
\[
\bigvee_{u \in A(v)} s(u) = \text{UNREACHABLE}
\quad \text{and} \quad
\not\exists h \in H : h \in A(v) .
\]
That is, every node in the ancestry of $v$ is either unreachable or terminates
before reaching a human anchor.
\end{definition}

\begin{定理}[Drift Probability Lower Bound]
Let $v$ be an agent at depth $d$ with ancestry $A(v) = \{v_0, v_1, \dots, v_d\}$
where $v_0 \in H$ and $(v_{i-1}, v_i) \in E$ for $i = 1, \dots, d$. Assuming each
edge fails independently with probability $p_f$ per time unit, the probability
that $v$ is drifted at time $t$ is at least

\[
P(\text{drift at } t) \geq 1 - (1 - p_f)^{d} .
\]

\textbf{Proof.} The agent is drifted if at least one edge in its ancestry chain
fails. The complement event—no edge fails—is $(1 - p_f)^d$. Taking the
complement yields the bound. $\square$
\end{定理}

This bound is independent of depth limits. A system with $d=5$ and $p_f = 0.01$
per hour has $P(\text{drift}) \geq 1 - (0.99)^5 \approx 0.049$ after one hour.
Depth limits do not change $p_f$; they only bound $d$, which appears in the
exponent, meaning diminishing $d$ reduces drift probability only logarithmically.

\begin{conjecture}[Drift Chain Growth]
If a drifted agent $v$ spawns children at rate $\lambda$ and each child is also
drifted, then the expected number of drifted agents $D(t)$ at time $t$ satisfies

\[
D(t) \geq D(0) \cdot e^{\lambda t} ,
\]

i.e., exponential growth in the number of drifted agents until external
intervention or resource exhaustion occurs.
\end{conjecture}

This model establishes that drift is not an edge case: under routine failure
conditions, it occurs with non-negligible probability even in small spawning
trees, and once it occurs, it grows without bound in the absence of external
intervention. Depth limits do not break the drift dynamics; they only limit the
initial depth of the tree.